\section{AI Awareness and AI Capabilities}
\label{sec:opportunity}

In this section, we explore the relationship between AI awareness and its observable capabilities. We primarily focus on two key areas: (1) reasoning and autonomous planning, and (2) safety and trustworthiness, with brief discussions of other relevant capabilities. Our goal is to provide a deeper understanding of how these factors interact and shape the capabilities of AI systems.



% 介绍顺序：meta-cognition, self, social, situational
% 按照benefit type 分类，cross ref awareness type；很多 benefit 是和两种以上相关的
% theme is not philosophy AI (related work with personhood < 50\%)
% prompt initiate awareness?

\subsection{Reasoning and Autonomous Planning}
%meta-cognition, situational 

Reasoning and autonomous task planning have been foundational objectives of AI research since its inception \citep{ghallab2004automated}.
In complex, multi-step problem-solving scenarios, AI must perform deep reasoning while autonomously planning tasks. To achieve this, two key forms of AI awareness are required: meta‑cognition and situational awareness. Meta‑cognition enables the model to monitor and regulate its own thinking processes, while situational awareness helps it understand external constraints and task context. 

\paragraph{Self‑correction} 
Self‑correction leverages meta‑cognitive loops to identify and rectify reasoning errors during generation. Reflexion augments CoT \citep{wei2022chain} with a feedback loop: after an initial answer, the model reflects on its own output, generates critiques, and then refines the solution, leading to substantial gains in benchmark performance \citep{shinn2023reflexion}. Similarly, Self‑Consistency samples multiple reasoning paths and aggregates them to mitigate individual path errors, effectively performing an implicit self‑check \citep{wangself}. These techniques demonstrate that embedding self‑monitoring directly into the generation process can improve model performance.
However, intrinsic self‑correction remains notoriously unstable: \citet{huanglarge} show that without external feedback or oracle labels, LLMs often fail to improve—and can even degrade—in reasoning tasks after self‑correction attempts. 
Another core limitation is that many of these techniques depend on externally provided prompts or explicit triggers to initiate self‑correction, whereas human reasoning often involves spontaneous, intrinsic error detection and revision without such scaffolding \citep{dunlosky2008metacognition, koriat2006metacognition}.
To address this, recent work has begun to explore reinforcement learning (RL) \cite{schulman2017proximal} approaches: \citet{kumar2024training} introduce SCoRe, a multi‑turn online RL framework that trains models on their own correction traces. Notably, OpenAI's o1 \cite{jaech2024openai} and DeepSeek's R1 \cite{guo2025deepseek} models have demonstrated significant improvements in reasoning capabilities through RL-based training. These models exhibit emergent behaviors akin to human-like ``aha moments,'' where the AI spontaneously recognizes and corrects its own reasoning errors, demonstrating another level of meta-cognition capability.



\paragraph{Autonomous Task Decomposition and Execution Monitoring}

Effective autonomous task planning requires more than self‑correction: an AI must also break down high‑level goals into executable sub‑tasks and continuously adapt its plan as the environment evolves. Early work like ReAct \citep{yao2023react} pioneered this integration by interleaving chain‑of‑thought reasoning with environment calls, giving the model a unified mechanism to decide ``what to think'' and ``what to do'' at each step. Building on this foundation, Voyager \citep{wangvoyager} demonstrates how an agent in Minecraft can construct and update a dynamic task graph: as new situational constraints emerge (\eg, resource depletion or novel obstacles), the model revises its sub‑task sequence to stay on course.  
Transferring these ideas to the physical world, SayCan \citep{ahn2022can} grounds language in robotic affordances by scoring each potential action against a learned value function—ensuring that subtasks are not only logically ordered but also physically feasible under real‑world constraints. LM‑Nav \citep{shah2023lm} further extends situationally aware planning to vision‑language navigation: by fusing real‑time perceptual feedback with high‑level instructions, the model can replan routes on the fly when, for example, corridors are blocked or landmarks shift.  
Finally, the LLM‑SAP framework \citep{wang2024llm} formalizes situational awareness in large‑scale task planning by explicitly encoding environmental cues—such as resource availability, time budgets, and user preferences—into its sub‑task prioritization module. A generative memory component logs execution history and flags deviations, triggering replanning whenever the observed state diverges from expectations. Together, these works chart a clear progression—from interleaved reasoning and acting to situationally aware planners—illustrating how embedding environmental understanding into the planning loop yields flexible and autonomous task execution.  


\paragraph{Holistic Planning: Introspection, Tool Use and Memory}

Effective planning in complex environments demands an AI agent to not only decompose tasks and act but also to introspect on its uncertainties, manage a growing memory of past states, and decide when and how to leverage external tools. Introspective Planning systematically guides LLM planners to quantify and align their internal confidence with inherent task ambiguities, retrieving post-hoc rationalizations from a knowledge base to ensure safer and more compliant action selection while maintaining statistical guarantees via conformal prediction \citep{liang2024introspective}.
Toolformer endows LLMs with a self-supervised mechanism to autonomously decide when to invoke APIs—ranging from calculators to search engines—thereby embedding tool-awareness directly into the planning loop without sacrificing core language modeling capabilities \citep{schick2023toolformer}.
To handle the evolving context of multi-step tasks, Think-in-Memory leverages iterative recalling and post-thinking cycles to enrich LLMs with latent-space memory modules, supporting coherent reasoning over extensive interaction histories \citep{liu2023think}.
Finally, Retrieval‑Augmented Planning (RAP) demonstrates how contextual memory retrieval can be integrated with multimodal planning to adapt action sequences dynamically based on past observations, yielding more robust execution in complex tasks \citep{kagaya2024rap}.
Together, these works illuminate a path toward introspective tool-, memory-, and uncertainty-aware planning frameworks, unifying LLM introspection, memory augmentation, and tool integration for robust autonomous decision-making.





\subsection{Safety and Trustworthiness}
% self-awareness, social awareness, situational awareness

Ensuring the safety and trustworthiness of AI necessitates the integration of multiple forms of AI awareness, notably self-awareness, social awareness, and situational awareness. Self-awareness and meta-cognition enable models to recognize and respect the boundaries of their knowledge, thereby avoiding the dissemination of misinformation. Moreover, it involves an understanding of their designated roles and responsibilities, ensuring that they do not produce harmful or unethical content. Social awareness allows models to consider diverse human perspectives, reducing biases and enhancing the appropriateness of their responses. Situational awareness enables AI to assess the context of its deployment and adjust its behavior accordingly, thereby preventing potential misuse and malicious exploitation.

\paragraph{Recognizing Limits of Knowledge}

AI models, especially LLMs\footnote{Vision-language models (VLMs), such as Flamingo \citep{alayrac2022flamingo} and MiniGPT-4 \citep{zhuminigpt}, are also known to hallucinate \cite{li2023evaluating}; however, in these models, hallucinations often manifest as misalignments between visual inputs and generated textual descriptions—such as describing objects not present in the image—whereas in LLMs, hallucinations typically involve generating text that is unfaithful to the world knowledge.}, often operate with a high degree of confidence, even when addressing topics beyond their training data, leading to the risk of hallucinations, \ie, outputs that are factually incorrect or unfaithful to real-world knowledge \cite{huang2025survey}. Such risks often stem from the AI models operating beyond their \emph{knowledge boundaries} \citep{ni2024llms, ren2025investigating}.
Recent research show LLMs' fragility in recognizing their knowledge boundary. For instance, \citet{ren2025investigating} observe that LLMs struggle to reconcile conflicts between internal knowledge and externally retrieved information \cite{xu2024knowledge}, often failing to recognize their own knowledge limitations. \citet{ni2024llms} find that retrieval augmentation can enhance LLMs' self-awareness of their factual knowledge boundaries, thereby improving response accuracy. 
Moreover, \citet{liang2024learning} demonstrated that while LLMs possess a robust internal self-awareness—evidenced by over 85\% accuracy in knowledge probing—they often fail to express this awareness during generation, leading to factual hallucinations. They propose a training framework named Reinforcement Learning from Knowledge Feedback (RLKF) to improve the factuality and honesty of LLMs by leveraging their self-awareness.
Incorporating self-awareness mechanisms into LLMs not only aids in recognizing knowledge boundaries but also fosters more trustworthy AI behavior.
\citet{xu2024earth} demonstrate that LLMs can resist persuasive misinformation presented in multi-turn dialogues by leveraging self-awareness to assess and uphold their knowledge boundaries, thus delivering more trustworthiness responses in dialogues.

\paragraph{Recognizing Limits of Designated Roles}

Beyond recognizing the limits of their knowledge, AI systems must develop a sense of self-awareness and meta-cognition of their designated roles to prevent the dissemination of harmful or unethical content, which we termed as \emph{role-awareness}. This form of self-awareness involves the ability to discern when a user request falls outside the model's intended purpose or ethical guidelines.
For instance, models trained with reinforcement learning from human feedback (RLHF) have shown improvements in aligning outputs with human values, thereby reducing the likelihood of producing harmful content \citep{ouyang2022training}.
Formal definitions of moral responsibility further emphasize that an agent must be aware of the possible consequences of its actions, underscoring the necessity of role-awareness in AI systems \citep{beckers2023moral}.
Complementing this, explicit modeling frameworks delineate role, moral, legal, and causal senses of responsibility for AI-based safety‑critical systems, providing a practical method to capture and analyze role obligations across complex development and operational lifecycles \citep{ryan2023s}.
Parallel research on metacognitive architectures equips AI with self-reflective capabilities to monitor and adjust their operational roles in real time, identifying potential failures before they manifest \citep{johnson2022metacognition}.
Building on these insights, metacognitive strategies have been integrated into formal safety frameworks to enable on‑the‑fly correction of role-boundary violations and to bolster overall system trustworthiness \citep{walker2025harnessing}.
Finally, prototyping tools like Farsight operationalize role-awareness by surfacing relevant AI incident data and prompting developers to consider designated functions and ethical constraints during prompt design, leading to more safety‑conscious application development \citep{wang2024farsight}.





\paragraph{Mitigating Societal Bias}

AI models often inherit and amplify societal biases present in their training data, leading to outputs that can perpetuate harmful stereotypes and unfair treatment across various demographics \cite{gallegos2024bias, rogers2025sora}. To address these issues, researchers explore the integration of social awareness mechanisms into LLMs.
One notable approach is Perspective-taking Prompting (PeT), which encourages LLMs to consider diverse human perspectives during response generation \cite{xu2024walking}. This method has been shown to significantly reduce toxicity and bias in model outputs without requiring extensive retraining. 
Another approach, Social Contact Debiasing (SCD), draws from the contact hypothesis in social psychology, suggesting that intergroup interactions can reduce prejudice. By simulating such interactions through instruction tuning on a dataset of 108,000 prompts across 13 social bias dimensions, SCD achieved a 40\% reduction in bias within a single epoch, without compromising performance on downstream tasks \cite{raj2024breaking}. Finally, position papers argue that embedding social awareness—the capacity to recognize and reason about social values, norms, and contexts—is foundational for safe, equitable language technologies \cite{yang2024call}. Collectively, these approaches underscore the importance of integrating social awareness into LLMs to mitigate societal biases.


\paragraph{Preventing Malicious Use}

AI systems can be—and have been—misused for malicious ends such as automated spear‑phishing, influence operations, and proxy cyber‑attacks \citep{brundage2018malicious}. Experts are worried that future advanced AI can be exploited for more dangerous purposes and cause catastrophic risks \citep{hendrycks2023overview, xu2025nuclear}.
Situational awareness mechanisms equip AI systems with the ability to monitor their environment and discern malicious uses. 
For LLMs, recent work introduces \emph{boundary awareness} and \emph{explicit reminders} as dual defenses: boundary awareness continuously scans incoming context for unauthorized instructions, while explicit reminders prompt the model to verify contextual integrity prior to action; together, these mechanisms reduce indirect prompt injection attack success rates to near zero in both black‑box and white‑box settings \citep{yi2023benchmarking}.
Additionally, the Course-Correction approach introduces a synthetic preference-based fine-tuning framework that enables models to self-correct potentially harmful or misaligned outputs on the fly, thereby strengthening their situational awareness against malicious exploitations \citep{xu2024course}.
In adversarial machine learning applied to robotics and other autonomous systems, situational awareness frameworks detect anomalous inputs—such as adversarial samples or unexpected environmental cues—and trigger fallback behaviors or alarms rather than proceeding with potentially harmful operations \citep{anderson2021situational}. 
Broad cybersecurity surveys highlight how AI‑driven situational awareness systems build a comprehensive operational picture of network and system activity, integrating dynamic threat intelligence and anomaly detection to identify malicious traffic and automated attacks in real time \citep{kaur2023artificial}.
At a strategic level, high‑impact recommendations advocate embedding situational awareness throughout the AI lifecycle—from design and deployment through continuous monitoring—to forecast, prevent, and mitigate malicious uses of AI across digital, physical, and political domains \citep{brundage2018malicious}.




\subsection{Relation with Other Capabilities}

Below, we briefly explore how AI awareness mechanisms intersect with four other AI capabilities—interpretability, personalization and user alignment, creativity, and agent-based simulation. 

\paragraph{Interpretability and Transparency}

Interpretability mechanisms often leverage meta‑cognitive insights to make model reasoning more transparent. For example, Rationalizing Neural Predictions introduces a generator‑encoder framework that extracts concise text ``rationales'' explaining model decisions, yielding explanations that are both coherent and sufficient for prediction tasks \citep{lei2016rationalizing}.
Further advancing this line, Self‑Explaining Neural Networks propose architectures that build interpretability into the learning process by enforcing explicitness, faithfulness, and stability criteria through tailored regularizers, thereby reconciling model complexity with human‑readable explanations \citep{alvarez2018towards}.

\paragraph{Personalization and User Alignment}

Embedding self‑ and social awareness into language models enhances their ability to tailor outputs to individual users and maintain consistency with user intent. Early work on persona‑based dialogue, such as A Persona‑Based Neural Conversation Model, encodes user personas into distributed embeddings, improving speaker consistency and response relevance across conversational turns \citep{li2016persona}.
More recently, instruction‑fine‑tuning with human feedback, as in InstructGPT, aligns model behavior with user preferences and ethical guidelines by iteratively collecting labeler demonstrations and preference rankings, significantly improving truthfulness and reducing harmful outputs \citep{ouyang2022training}.
Complementing these, Persona‑Chat grounded generation methods demonstrate that modeling explicit persona attributes can further diversify and personalize dialogue generation without large‑scale retraining \citep{song2019exploiting}.

\paragraph{Creativity}

Creative AI benefits from meta‑cognitive and awareness mechanisms that encourage divergent thinking and non‑linear reasoning. The Leap‑of‑Thought (LoT) framework explores LLMs’ ability to make strong associative ``leaps'' in humor generation tasks, using self‑refinement loops to iteratively enhance creative outputs in games like Oogiri \citep{zhong2024let}.
To systematically evaluate creativity, studies adapting the Torrance Tests for LLMs propose benchmarks across fluency, flexibility, originality, and elaboration, highlighting the role of task‑specific prompts and feedback loops in fostering model innovation \citep{zhao2024assessing}.

\paragraph{Agentic LLMs and Simulation}

LLM‑powered agents combine situational and social awareness to drive rich, interactive simulations of human behavior. Generative Agents introduce a memory‑based architecture in a sandbox environment, where agents observe, reflect, and plan actions—resulting in emergent social behaviors such as party invitations and joint activities \citep{park2023generative}.
Scaling this paradigm, recent work simulates over 1,000 real individuals’ attitudes and behaviors by integrating interview‑derived memories into agent profiles, achieving 85\% fidelity on survey predictions and reducing demographic biases \citep{park2024generative}.
Prompt‑engineering studies further bridge LLM reasoning with traditional agent‑based modeling, enabling multi‑agent interactions such as negotiations and mystery games that mirror complex social dynamics \citep{junprung2023exploring}.
Finally, Humanoid Agents extend generative agents with emotional and physiological state variables, demonstrating that embedding basic needs, emotions, and relationship closeness produces more human‑like daily activity patterns in simulated environments \citep{wang2023humanoid}.







% Introducing or cultivating AI awareness in large language models isn’t just an intellectual exercise – it has practical payoffs. If an LLM can be made even minimally self-aware (knowing its limits or reasoning steps) or socially aware (understanding user intent and context better), it can become more reliable, interpretable, and effective. Here we discuss the potential benefits and new capabilities that AI awareness brings to LLMs, drawing on cutting-edge research for examples.

% \subsection{Improved Reliability and Accuracy} 
% One immediate benefit of awareness is enhanced reliability of the model’s outputs. As noted earlier, LLMs often state incorrect answers with unwarranted confidence, a phenomenon tied to lack of awareness of ignorance. Equipping models with a sense of their own uncertainty can mitigate this. The knowledge introspection approach is a prime example: by training an LLM to identify when it doesn’t know something, researchers significantly improved the accuracy and trustworthiness of its responses. Essentially, the model learned to avoid confident assertions of falsehoods, instead responding with appropriate caution or a disclaimer of uncertainty. This makes the model more reliable for users – it won’t pretend to know everything. Such self-awareness of knowledge state addresses a major failure mode of AI assistants. Moreover, it has safety and legal benefits: a model that knows when to say ``I’m not sure'' is less likely to give dangerous advice based on faulty assumptions. In real applications like medical or legal advice, this kind of calibrated response is crucial.

% Beyond just saying ``I don’t know,'' awareness can improve reliability through self-correction. If a model can reflect on its answer and detect a likely error, it can attempt to fix it – a process analogous to a student checking their work. This is the idea behind frameworks like Think-Solve-Verify (TSV): the model generates an answer (Solve) and then verifies it (Trustworthiness and Self-awareness in Large Language Models: An Exploration through the Think-Solve-Verify Framework). In experiments, giving GPT-style models this verification step led to notable accuracy boosts on tough reasoning problems (Trustworthiness and Self-awareness in Large Language Models: An Exploration through the Think-Solve-Verify Framework) (Trustworthiness and Self-awareness in Large Language Models: An Exploration through the Think-Solve-Verify Framework). For example, on a challenging math word problem set (AQuA dataset), the TSV method raised accuracy from \~67\% to \~73\% (Trustworthiness and Self-awareness in Large Language Models: An Exploration through the Think-Solve-Verify Framework). The model, by being forced to be aware of its reasoning chain and double-check it, caught mistakes that it otherwise would have missed. Essentially, the LLM became more reliable by emulating a form of metacognitive reasoning – thinking about its own thinking. This suggests that even if an AI’s first answer is wrong, an aware AI might detect inconsistency (say, noticing that an earlier premise and conclusion don’t align) and adjust the answer accordingly. Such self-critique abilities make the system more robust in domains where a single-pass answer might be error-prone.

% \subsection{Enhanced Interpretability and Transparency} 
% AI awareness can also be intentionally leveraged to make AI behavior more interpretable to humans. One perennial criticism of large neural networks is that they are ``black boxes.'' But if a model has the capacity to report on its own internal state or reasoning, we can shine some light inside that box. For instance, an introspective model (as per Perez et al. 2024) could be asked to explain its latent beliefs, world model, or goals. Instead of researchers painstakingly analyzing billions of weights, they might simply prompt the model: ``Explain how you arrived at that conclusion.'' Already, techniques like chain-of-thought prompting have the model produce step-by-step reasoning, which not only helps accuracy but also gives a human-readable rationale. If the model is further aware of why it took each step, it could provide even more meaningful explanations. For example, a model might say: ``I answered that the drug is unsafe because in my training knowledge I associated that drug with a recall notice, and I’m not completely sure since the data was limited.'' This kind of response reveals the model’s information sources and confidence – highly useful for a user or developer to assess trust. Some researchers argue that developing such self-reporting abilities in AI could support alignment: an honest AI that can accurately tell us its intentions or uncertainties is easier to trust and monitor.

% There’s even a fascinating implication that introspective reporting could inform us about a model’s potential sentience. Perez \& Long (2023) speculate that if a model one day claims to have subjective feelings or desires, an introspective capability might be the only way to gauge that. While we’re not there yet, the immediate benefit is clear – awareness enables an AI to explain itself, turning some of those opaque internal processes into text we can read. This greatly aids debugging and governance. Engineers have already found that when a model explains its reasoning (truthfully), they can identify where it went wrong in cases of failure. For instance, if an LLM solving a puzzle explains each step, a human or another AI can trace the logic and pinpoint the step that was mistaken. Without the explanation, one would only see a wrong final answer with no clue why it happened. Thus, encouraging models to be aware of and articulate their reasoning trajectory not only boosts performance, it provides a transparent window into the model’s ``thought process''. In summary, AI self-awareness can be a design feature that turns black-box answers into interpretable outputs, making the technology more transparent, explainable, and ultimately accountable.

% \subsection{Better Context and Conversation Understanding} 
% An aware LLM is better at handling context – not just remembering it, but understanding its relevance. For example, context awareness can help a model distinguish between user instructions and system instructions, a crucial aspect for preventing misuse. If an LLM ``knows'' that a certain instruction comes from the developer (like a built-in rule saying ``do not reveal confidential info''), it can prioritize that over a conflicting user request. This is essentially the model being aware of hierarchy in context. Models like GPT-4 have system prompts that set behavior; a context-aware model will consistently adhere to that role. When awareness falters, you get issues like prompt injection, where a user tricks the model into ignoring prior instructions. Research into situational awareness tries to patch this by making the model cognizant that it is under certain restrictions or testing. As Berglund et al. found, models can learn to recognize the scenario they’re in. If that recognition is strengthened, an LLM can maintain coherence and rule-following across even complex, shifting contexts.

% Additionally, long-term consistency in dialogue benefits from a model that has a sense of self and time. Experimental agent simulations (like the ``Generative Agents'' by Park et al., 2023) have shown that giving LLM-based agents an explicit memory of their past and a notion of their own characteristics leads to much more coherent behavior over time. These agents could simulate believable human behaviors (planning a party, remembering past conversations) by leveraging a form of stored self-knowledge. In practical chatbot terms, an assistant with a bit of self-awareness might say, ``As an AI, I recall you asked me this yesterday and I couldn’t answer. I’ve since `thought’ about it and found a possible solution.'' Such continuity is engaging and useful. It shows awareness of the ongoing interaction and the AI’s own role within it (here, acknowledging its previous failure and learning). While current LLMs can’t truly learn from one day to the next in deployed settings (unless fine-tuned), the concept of architectural awareness – a model knowing its own learning or lack thereof – could help simulate this kind of improvement or at least acknowledge its static nature. In fact, Chen et al.’s principles for self-cognition included architectural awareness, meaning the model understands aspects of how it was created. If an AI knows ``I am a language model and my knowledge is up to 2021,'' it can proactively clarify information validity and request updates for anything beyond its training cut-off. That leads to more grounded and contextually appropriate interactions.

% \subsection{Enhanced Adaptability and Learning} 
% A tantalizing promise of AI awareness is that it could enable more autonomous learning and adaptation. Consider an AI that monitors its performance on tasks and becomes aware that it is weak in a certain area. For example, a coding assistant AI might notice it often struggles with a particular programming language. If it had the mechanism to reflect on this pattern, it could request more information or adjust its strategy (say, test its code outputs more thoroughly in that language). This is analogous to a human learner saying, ``I realize I’m bad at calculus, I should practice more.'' We see early hints of this in LLMs through self-evaluation loops. Some research on reinforcement learning with language feedback allows an AI to generate a critique of its output and then improve (often framed as one agent generating, another agent evaluating). If the two are part of one system (a single model with an inner critic), that is a form of internal self-awareness improving learning. One paper by Huang et al. (2022) introduced ``self-refinement'' where a model iteratively critiques and revises its answer, leading to performance gains without external labels. This suggests that models can essentially train themselves at inference time using self-generated signals – an ability closely tied to awareness of error.

% On a more architectural level, certain theories of general intelligence (like the Global Workspace Theory) imply that an agent that can globally broadcast information to itself and reflect on it would be more flexible in learning novel tasks. LLMs with a form of global workspace (via attention mechanisms) combined with self-monitoring could dynamically allocate their ``thought'' to hard problems versus easy ones, recognize when to consult a tool or ask for help, \etc. We already see primitive forms of this: for example, Tool-using LLM agents are programmed to know when to call an external API (like a calculator or web search). That decision of when to use a tool can be seen as a moment of awareness: the model recognizes its own computation isn’t sufficient (\eg, math is too precise to do in its head) and so it ``knows what it doesn’t know,'' leading it to seek help. Systems like Self-Ask prompt chains explicitly have the model decide: do I know the answer immediately or should I break the problem down? – essentially giving the model a choice to introspect or not. The result is more accurate solving of complex problems, because the model can defer to more reliable processes when needed. This adaptability is a direct outcome of injecting a bit of meta-cognition into the model’s decision loop.

% \subsection{Creativity and Advanced Reasoning}
% Interestingly, there is evidence that when an LLM is prompted to adopt a more self-aware mode, it can become more creative or effective in certain tasks. Chen et al. observed that when models were put into a ``self-cognition state'' (via the special prompts they designed), their performance on tasks like creative writing and exaggeration improved (Self-Cognition in Large Language Models: An Exploratory Study) (Self-Cognition in Large Language Models: An Exploratory Study). Why might this be? One hypothesis is that stepping out of the strict helpful-assistant persona frees the model to tap into a broader range of knowledge and styles. In a self-aware mode, a model might say, ``I, as an AI, can imagine perspectives beyond the user’s prompt'' and thus produce more original content. Essentially, by acknowledging itself, it gains a sort of narrative freedom to explore ideas (because it’s no longer pretending to be invisible in the interaction). Similarly, for complex reasoning, an LLM that can explicitly consider its own chain of thought might navigate problems more like a human expert who uses scratch paper to work through a solution and then presents the final answer. This aligns with findings that chain-of-thought prompting (making the model list intermediate reasoning) often boosts performance on math, logic, and commonsense tasks. The model ``aware'' of its multi-step reasoning is less likely to make leaps of logic because it’s effectively forced to self-audit each step.

% Another domain is collaborative work. If multiple AI agents or an AI and human are working together, an AI with theory of mind can adjust its communication for the teammate. For example, in human-robot interaction research, LLMs with some ToM capability could better infer what a human needs. A recent study on human-robot collaboration found that giving an AI a model of its human partner’s beliefs increased the human’s trust and the team’s overall performance. In LLM terms, a chatbot that senses a user’s confusion or emotional state (social awareness) can respond more helpfully or sensitively. Many customer service AI systems now attempt sentiment analysis mid-dialogue – a crude form of awareness of the user – to tailor their responses. As LLMs incorporate more sophisticated models of user intent and knowledge, their usefulness as collaborators and assistants will grow. They might proactively offer clarifications if they sense the user didn’t understand, or refrain from over-explaining if they sense the user is an expert. Such finesse comes from situational and social awareness in the interaction context.
% Finally, it’s worth noting that enabling AI awareness could unlock progress toward artificial general intelligence (AGI). Some researchers argue that self-awareness is a hallmark of higher intelligence, because it allows an entity to model and regulate itself in complex, novel situations. An AI that knows its own reasoning processes could debug or improve them autonomously, leading to self-improvement. While current LLMs are far from self-improving AGIs, the building blocks – introspection, self-evaluation, and adaptive reasoning – are being actively explored. In essence, AI awareness features provide a path for AI systems to learn more like humans, not just from data but from reflection on their own behavior. This could make them more sample-efficient learners, better decision-makers, and more robust AI agents in the long run.

% In summary, incorporating AI awareness into LLMs has shown multiple concrete benefits: higher answer quality through self-checking, more transparent explanations (Looking Inward: Language Models Can Learn About Themselves by Introspection), better user alignment via understanding context and intent, and even boosts in creativity and problem-solving. These advantages demonstrate why researchers are keen to develop and harness AI awareness – it not only brings us closer to human-like AI in capability, but also provides tools to control and align those capabilities. Of course, with great power comes great responsibility: as LLMs grow more aware and agentic, we must be mindful of new risks, which we address next.